\section{引言}

一纸传单能让一个王朝的朝堂血流成河，这在今天听起来像是修辞。万历三十一年（1603）的「妖书案」却是这样的事件：
一份仅几百字的匿名文字在京师传开，短短十几天内朝廷人人皆知；它没有议论、没有叙述，只有两个人的对话，
却把立储、阁臣更替、后宫干政全串成了一条线索。本文要问的不是这桩公案谁是真凶，而是一个可计算的问题：
**在没有报纸与电报的时代，一部通俗史著把帝国的「信息战」写到了什么颗粒度，又把一个什么样的信息结构
交给了知识图谱。**

本文把全书 156 章抽取的明史知识图谱（1231 人物、563 地点、1106 事件、2324 条关系）当作观测对象，
只取其中与信息生产、传播、封锁有关的那一层，回答三件事。

\textbf{贡献一：给出信息型关系的可抽取口径。}全书 2324 条关系归入 12 类，「政治攻讦」193 条（8.3\%）。
在这一类关系的文本里，含弹劾、告发、诬告、构陷、监视等**信息型动作**的占 24.4\%，
同为攻击性的「敌对冲突」仅 11.2\%，纯从属的「君臣从属」仅 0.8\%。把攻击按「是否经由信息渠道」切开，
是既有抽取体系里原本没有的一个维度。

\textbf{贡献二：量化攻讦关系随叙事推进的上升。}以书的七部为时间代理，政治攻讦占该部全部关系的比重
从第 1--2 部的 6.4\% 升到第 6--7 部的 16.0\%（$2.48$ 倍）；对照组「敌对冲突」同期从 22.7\% 变到 22.1\%
（$0.98$ 倍，基本不动）。信息型内耗的上升不是军事对抗上升的副产品。

\textbf{贡献三：信息型机构在图谱里是隐形的。}角色文本含锦衣卫、东厂、西厂的 29 名人物，平均关系度数 3.76，
与全库均值 3.78 几乎相同（比值 $0.99$），却只占攻讦关系端点的 3.6\%；而同代这套制度的内廷宦官 51 人，
平均度数 6.10（$1.61$ 倍），占攻讦端点 10.1\%。书里写得最多的是机构，进到图谱结构里最深的是它的代言人。

边界必须先写清楚。本文用「信息战」一词是**现代类比**，不是对明代的史实定性：明代没有近代新闻业，
也没有「信息战」这个题名，本文只在其可比层面使用这个名词。同时，通俗史著对厂卫的描写带有明显的文学化成分，
其可信度低于档案，因此本文所有涉及机构行为的判断都只落在**本书语料内**。
这篇论文不主张发现了明代的新史实，只主张：一部通俗史著的信息结构可以被抽取，而且抽出来的结构
与作者的笔法之间存在可测的错位。
